% Standalone problem document, generated from the adjacent README.md.
% Compile with XeLaTeX. Mathematical content is maintained in Markdown.
\documentclass[11pt,a4paper]{article}
\usepackage[fontsize=10.5pt]{scrextend}
\usepackage[margin=22mm,headheight=15pt,headsep=9mm,footskip=11mm]{geometry}
\usepackage{fontspec}
\setmainfont{texgyrepagella}[Extension=.otf,UprightFont=*-regular,BoldFont=*-bold,ItalicFont=*-italic,BoldItalicFont=*-bolditalic]
\setsansfont{texgyreheros}[Extension=.otf,UprightFont=*-regular,BoldFont=*-bold,ItalicFont=*-italic,BoldItalicFont=*-bolditalic]
\setmonofont{texgyrecursor}[Extension=.otf,UprightFont=*-regular,BoldFont=*-bold,ItalicFont=*-italic,BoldItalicFont=*-bolditalic,Scale=MatchLowercase]
\usepackage{amsmath,amssymb}
\usepackage{unicode-math}
\setmathfont{texgyrepagella-math.otf}
\usepackage{xcolor,microtype,enumitem,needspace,fancyhdr}
\usepackage{bookmark}
\definecolor{ink}{HTML}{17344B}
\definecolor{accent}{HTML}{197D80}
\definecolor{muted}{HTML}{596773}
\hypersetup{colorlinks=true,linkcolor=accent,urlcolor=accent,pdftitle={AC-07 - Scholz--Brauer
inequality for multiplication
chains},pdfauthor={Open Problems in Numerical Linear Algebra}}
\urlstyle{same}
\setlength{\parindent}{0pt}
\setlength{\parskip}{5pt plus 1pt minus 1pt}
\linespread{1.04}
\setlength{\emergencystretch}{3em}
\widowpenalty=10000
\clubpenalty=10000
\displaywidowpenalty=10000
\allowdisplaybreaks[1]
\setlist{nosep,leftmargin=1.5em}
\providecommand{\tightlist}{\setlength{\itemsep}{0pt}\setlength{\parskip}{0pt}}
\providecommand{\pandocbounded}[1]{#1}
\pagestyle{fancy}
\fancyhf{}
\fancyhead[L]{\sffamily\footnotesize\color{muted}OPEN PROBLEMS IN NUMERICAL LINEAR ALGEBRA}
\fancyhead[R]{\sffamily\bfseries\footnotesize\color{accent}AC-07}
\fancyfoot[L]{\parbox[b]{0.9\textwidth}{\sffamily\fontsize{8}{10}\selectfont\color{muted}Editorial ratings; a bounded search cannot certify that a problem remains unsolved.\\Literature check: 2026-09-08}}
\fancyfoot[R]{\sffamily\footnotesize\color{muted}\thepage}
\renewcommand{\headrulewidth}{0.3pt}
\renewcommand{\headrule}{\hbox to\headwidth{\color{accent}\leaders\hrule height \headrulewidth\hfill}}
\makeatletter
\renewcommand{\section}{\Needspace{5\baselineskip}\@startsection{section}{1}{0pt}{12pt plus 2pt minus 2pt}{4pt}{\sffamily\bfseries\large\color{ink}}}
\renewcommand{\subsection}{\Needspace{4\baselineskip}\@startsection{subsection}{2}{0pt}{9pt plus 1pt minus 1pt}{3pt}{\sffamily\bfseries\normalsize\color{ink}}}
\makeatother
\setcounter{secnumdepth}{0}
\begin{document}
{\sffamily\small\color{muted}Arithmetic and complexity\par}
\vspace{3pt}
{\raggedright\hyphenpenalty=10000\exhyphenpenalty=10000\sffamily\bfseries\fontsize{21}{24}\selectfont\color{ink}Scholz--Brauer
inequality for multiplication chains\par}
\vspace{7pt}
{\sffamily\small\textbf{Difficulty:} challenging\quad\textbf{Importance:} interesting
to specialist\par}
\vspace{4pt}
{\color{accent}\hrule height 0.8pt}
\vspace{6pt}
\textbf{Topic:} multiplication counts for matrix powers

\textbf{Status:} open; admitted after a bounded literature search found
no resolution.

\subsection{Problem statement}\label{problem-statement}

An addition chain for a positive integer \(n\) is a sequence
\(1=a_0<a_1<\cdots<a_r=n\) in which every \(a_i\) with \(i>0\) is a sum
of two earlier terms, which may coincide. Let \(\ell(n)\) be its minimum
possible length \(r\). Is

\[\ell(2^n-1)\le n+\ell(n)-1\qquad(n\ge1)?\]

The case \(n=1\) uses \(\ell(1)=0\). A chain implements a computation of
\(A^n\) by matrix multiplications; the question is about this
multiplication-only model, not every possible matrix-function algorithm.

\subsection{References and status}\label{references-and-status}

Bläser,
\href{https://theoryofcomputing.org/articles/gs005/gs005.pdf}{2013},
Research Problem 2.9. N. Clift,
\href{https://www.additionchains.com/ExactScholzBrauer.pdf}{\emph{A
Short Note on Exact Equality for the Scholz--Brauer Conjecture}} (2024),
distinguishes this inequality from the stronger equality and gives a
counterexample to equality. T. Agama,
\href{https://www.researchgate.net/publication/397490267_A_PROGRESS_ON_THE_SCHOLZ_CONJECTURE_ON_ADDITION_CHAINS}{\emph{A
Progress on the Scholz Conjecture on Addition Chains}}, manuscript dated
10 May 2026, §1, still states the general inequality as open; its result
imposes additional conditions on an optimal addition chain. Searches for
``Scholz Brauer conjecture proof 2026'' and inspection of these sources
found no proof or disproof of the unrestricted inequality.
\textbf{Admitted: no resolution located.}
\end{document}
